% Folder 146, batch 04 — archivists' pages 61–80.
% Pass: Opus 5.5 (claude-opus-5-5), 2026-10-02 — first pass, unchecked against the pages by a human.
% Skipped: 68 (unrelated typed letter), 69, 75, 79 (blank or bleed-through only), 71 (typed exam subject, nothing in his hand).
\documentclass[11pt,a4paper]{article}
\input{../preamble/grothendieck}

\folder{146}
\batch{4}
\pages{61}{80}
\foldertitle{Groupoïdes de Teichmüller et les S0 Tg,!ν, S0 TG : notes manuscrites (s.d.).}
\dating{[à partir de 1978]}
\watermark{Édition de démonstration}

\begin{document}

\page{61}
\drawing{En tête de page, plusieurs croquis : un octaèdre (ou double pyramide) aux faces hachurées ; un triangle avec un petit triangle intérieur ; deux triangles réunis par un sommet (« nœud papillon ») ; un triangle marqué $+$, avec au-dessous un segment et un triangle aux sommets $0$, $1$, $\infty$.}

\drawing{Dans la colonne de gauche, trois sphères : la première porte un cercle équatorial avec cinq points marqués ; la deuxième un sommet au pôle relié par des arcs à quatre points de l'équateur ; la troisième des cercles méridiens et un axe pointillé, avec des points marqués aux pôles et sur l'équateur.}

En regard des sphères :
\[
\mathbb{Z}/5\mathbb{Z} \;\to\; \mathbb{D}_5
\]
(\(\mathbb{D}_5 \times \mathbb{Z}/2\mathbb{Z}\) dans \uncertain{le groupe} non orienté) ;
\(\mathbb{Z}/2\mathbb{Z}\) (\(\mathbb{D}_2\) dans \uncertain{le groupe} non orienté) ;
\[
\mathbb{Z}/3\mathbb{Z} \;\to\; \mathfrak{S}_3
\]
(\(\mathfrak{S}_3 \times \mathbb{Z}/2\mathbb{Z}\) dans \uncertain{le groupe} non orienté).

12 possibilités \ill{} (\uncertain{triangul.} pentag. sur $K$)
\uncertain{CAD} \struck{\ill{}} 12 sommets de l'icos.
le stabilisateur \uncertain{de chacun} est
$\mathbb{Z}/5\mathbb{Z}$ (\add{\uncertain{dans le groupe non orienté}} \struck{\uncertain{diédral} $\mathbb{Z}/5\mathbb{Z}$}, $\mathbb{D}_5$)
\uncertain{Implique orientation de $K$}

$30 = 15\times 2$ possibilités (arêtes orientées \add{\ill{}} des $K$)
\uncertain{CAD} 30 arêtes de l'icosaèdre, le stabilisateur d'une arête est $\mathbb{Z}/2\mathbb{Z}$
(et $\mathbb{D}_2 = \mathbb{Z}/2 \times \mathbb{Z}/2$ dans \uncertain{non} orienté)
\uncertain{n'implique pas orientation de $K$}

$20 = 10\times 2$ possibilités (triangles de $K$ \struck{\ill{}} \add{orientés} de $K$)
(NB 20 faces de l'icosaèdre, stabilisateur de chacune est $\mathbb{Z}/3\mathbb{Z}$ ;
\uncertain{groupe} $\mathfrak{S}_3$ \uncertain{dans non} orienté)

\page{62}
Avant de me lancer vers une exploitation des principes dégagés, il me faut encore développer la formulation des « points base à l'infini ». \struck{\ill{}} (NB Le point de départ est le suivant. \struck{\ill{}} Le point de départ est la dimension 1, \struck{bien} \add{mais} \uncertain{pour une bonne partie} il sera nécessaire de revenir par la suite sur les phénomènes analogues en dimension quelconque, relatifs à la donnée d'un diviseur à croisements normaux…) Soit $X$ un schéma local régulier de dim.\ 1 (i.e.\ le spectre d'un anneau de valuation discrète \uncertain{hensélien}), $x$ son point fermé, $k = k(x)$,
\[
\mathcal{U} = X \setminus \{x\} = \{\eta\} \tag{53}
\]
\struck{$\mathcal{U}$ =} son point générique, et l'\uncertain{espace} \uncertain{vectoriel}
\[
T_x = \mathbb{V}(\mathfrak{m}_x/\mathfrak{m}_x^2) = \check{\mathbb{V}}\bigl((\mathfrak{m}_x/\mathfrak{m}_x^2)\bigr) \tag{54}
\]
le $k$-schéma « espace tangent en $x$ », \(T_x^* = T_x \setminus \{0\}\), \uncertain{qui est épointé}. On a alors un morphisme canonique
\[
\hat{\Pi}_1(T_x^*) \longrightarrow \hat{\Pi}_1(\mathcal{U}) \tag{55}
\]
(où les $\hat{\Pi}_1$ ici \uncertain{sont} les $\Pi_1$ schématiques \add{\uncertain{\ill{}}}). \struck{On sait \ill{} \ill{}} (\uncertain{puisque} $\hat{\mathcal{U}} \to \mathcal{U}$ \uncertain{donne un} \uncertain{iso} $\hat{\Pi}_1(\hat{\mathcal{U}}) \to \hat{\Pi}_1(\mathcal{U})$) \struck{\ill{} sait alors} \uncertain{complété} \ill{}. Alors on sait qu'on a

\note{La page s'arrête sur « Alors on sait qu'on a » ; la suite ne se trouve pas dans ce lot. Les numéros (53)--(55) prolongent une numérotation commencée avant le lot.}

\page{63}
\note{Page presque entièrement faite de figures ; seuls les quelques mots et symboles qui les accompagnent sont transcrits.}

\drawing{En haut à gauche, un polyèdre en perspective formé de deux faces pentagonales (arêtes pleines et pointillées), les arêtes marquées $a, b, c, d, e$ et $a', b', c', d', e'$ ; une flèche courbe près de $b'$.}

pentagone orienté \qquad $b'c'd'e'$\note{Lecture douteuse : on pourrait lire aussi « $b'e'd'e'$ ».}

\drawing{Une série de sphères portant des points cerclés sur l'équateur et aux pôles, reliés par des arcs méridiens, étiquetées : $D_{01}$ ; $D_{12}$ (deux sphères) ; $D_0$ (cinq points sur un cercle équatorial, une flèche, signe $+$) ; $D_1$ (un pôle relié à quatre points de l'équateur, avec la mention « $(+\circlearrowleft)$ ») ; \struck{$D_2$} $D_2$ (sphère marquée $+$ et $-$).}

\drawing{Dans la moitié inférieure : un petit triangle hachuré ; un icosaèdre en perspective dont les arêtes portent les lettres $a, b, c, d, e$ ; un grand patron déplié de faces (pentagones et quadrilatères) aux arêtes marquées $a, \ldots, e$, avec des flèches courbes indiquant des recollements ; à droite, trois sphères au crayon portant des points cerclés et des méridiens.}

\page{64}
\ill{} les \uncertain{étaient} $0, \infty$ au lieu de $0, 1$, et la \uncertain{position} (\uncertain{lieu} \ill{} de $0, \infty$ sur $\mathbb{X}_1$) était $1$ au lieu de $\infty$, \struck{i.e.\ \ill{} \ill{} --- $\mathbb{X}_1$ \ill{} \ill{} but et \ill{} $\mathbb{X}_0$}, l'application \uncertain{dont il} serait donné par (25), \add{\uncertain{ou dans la composition avec une transformation homographique}}
$z \mapsto z^2$, \struck{\ill{} \ill{} \ill{} \uncertain{composée} avec l'isomorphisme}
\[
g_1(z) = \psi\bigl(\varphi(z)^2\bigr),
\]
\struck{(\uncertain{application} homographique \ill{} \ill{} qui réalise} où $\psi$ est l'application homographique qui réalise
\[
0 \;\; 1 \;\; \infty \longrightarrow 0 \;\; \infty \;\; 1 ,
\]
i.e.\ \struck{$\psi(z) = \sigma_0(z) = \frac{z}{z-1}$}, et où $\varphi$ est l'application homographique qui réalise \struck{donc on trouve}
\[
1 \mapsto 0, \quad -1 \mapsto \infty, \quad 0 \ (\text{ou } \infty) \mapsto 1 ,
\]
\struck{$\mathbb{X}_1 \to \mathbb{X}_0$, $z \mapsto \frac{z^2}{z^2-1}$} \quad $\varphi(z) = \dfrac{z-1}{z+1}$, donc
\[
\psi(z) = \sigma_0(z) = \frac{z}{z-1},
\]
et en composant l'application \ill{} :
\[
\left\{
\begin{aligned}
& g_1 : \mathbb{X}_1 \longrightarrow \mathbb{X}_0 : \\
& g_1(z) = -\frac{1}{4}\,\frac{(z-1)^2}{z}
\end{aligned}
\right. \tag{26}
\]
\note{Une première écriture, lue avec doute $z \mapsto z^4/(z^4-1)$, est biffée à droite de « $g_1 : \mathbb{X}_1 \to \mathbb{X}_0$ ».}
\marginal{équation du monogone, \ill{} !}

d'où comme équation pour le bigône
\[
\left\{
\begin{aligned}
& g_2 : \mathbb{X}_2 \longrightarrow \mathbb{X}_0 : \\
& g_2(z) = g_1(z^2) = -\frac{1}{4}\,\frac{(z^2-1)^2}{z^2} = -\frac{1}{4}\Bigl(z - \frac{1}{z}\Bigr)^2
\end{aligned}
\right. \tag{27}
\]
(équation du \emph{bigône}, \uncertain{de merveille} !)

\note{Le début de la page (« \ill{} les étaient $0, \infty$ au lieu de $0, 1$ ») et le renvoi à (25) supposent un argument commencé avant ce lot.}

\page{65}
\note{Les notes manuscrites occupent le haut d'une feuille dont l'imprimé, tête-bêche, est daté de Montpellier, 31 mai 1979 : c'est un terme a quo pour cette page. Seules les notes de Grothendieck sont transcrites.}

Un tableau, colonnes séparées par des traits ; en tête des deux dernières paires de colonnes : « \uncertain{NB des cas} » et « \uncertain{Nb de feuillets} ».
\[
\begin{array}{cc|c|c||cc|cc}
\mathbb{D}_5 & \mathbb{D}_5 \times \mathbb{Z}/2\mathbb{Z} & \mathbb{Z}/5\mathbb{Z} & \mathbb{D}_5 & 12 & 12 & 12 & 6 \\
\mathbb{Z}/2\mathbb{Z} & \mathbb{D}_2 \times \mathbb{Z}/2\mathbb{Z} & \mathbb{Z}/2\mathbb{Z} & \mathbb{Z}/2\mathbb{Z} \times \mathbb{Z}/2\mathbb{Z} & 30 & 15 & 30 & 15 \\
\mathbb{Z}/3\mathbb{Z} & \mathfrak{S}_3 \times \mathbb{Z}/2\mathbb{Z} & \mathbb{Z}/3\mathbb{Z} & \mathfrak{S}_3 & 20 & 10 & 20 & 10
\end{array}
\]
\note{Dans la première ligne, un nombre biffé et noirci précède le second « 12 ». Dans la première colonne, l'entrée de la troisième ligne est surchargée et entourée ; le « $\mathbb{Z}/2\mathbb{Z}$ » de la deuxième ligne pourrait se lire « $\mathbb{Z}/4\mathbb{Z}$ ».}

\drawing{Sous le tableau : un hexagone (deux triangles accolés et leurs diagonales) dont les côtés et les segments intérieurs portent les lettres $a, b, c$ et $a', b', c'$, avec une flèche de rotation ; un carré ; un pentagone, accompagné de la mention « cas d'incidence » ; une ellipse, accompagnée de la mention « pentagone orienté ».}

\page{66}
\note{Page de calculs et de figures écrite dans les deux sens : la moitié supérieure à l'endroit, la moitié inférieure tête-bêche. On transcrit d'abord la première, puis la seconde, qui porte dans son coin le nombre « 99 », peut-être une pagination de l'auteur.}

\(-1 \mapsto (-1)\) \;/\; \(0\) \;/\; \(1 \mapsto 0\)

\begin{tikzcd}[column sep=large]
X_2 \arrow[r, "z \mapsto z^2"] & X_1 \arrow[r] & X_0
\end{tikzcd}

\[
\begin{array}{ccc}
X_2 & X_1 & X_0 \\
0 & 0 \,|\, \infty & \infty \\
-1 & 1 & 0 \\
\infty & \text{\struck{$-1$}}\ \ast & 1
\end{array}
\]
\note{Le signe $\ast$ marque un symbole noirci illisible, écrit sur le $-1$ biffé. À gauche de la colonne de $X_2$, un « 1 » isolé en face de $-1$.}

\[
\frac{(1+z)^2 - (1-z)^2}{(1-z)^2} \Big/ \frac{(1+z)^2}{(1-z)^2}
= \frac{1 - \left(\frac{1+z}{1-z}\right)^2}{\left(\frac{1+z}{1-z}\right)^2},
\qquad = \frac{-4z}{(1-z)^2}
\]
\note{Calcul griffonné ; la disposition des exposants et des fractions est reconstituée avec doute.}

\[
-\frac{1}{4i}\,(i-1)^2 = \frac{1}{2}, \qquad (i-1)^2 = -1+1-2i .
\]

\drawing{Une ellipse dont un arc épais va d'un point marqué à $1$, avec $\infty$ en bas ; une sphère dont l'équateur porte un segment épais entre deux points marqués $0$, avec $1$ au milieu et $\infty$ en avant.}

\emph{Moitié tête-bêche.}

\drawing{Un quadrilatère de sommets $0, 1, \infty$ (et un quatrième), coupé par une diagonale, deux régions hachurées ; une ellipse traversée par le segment $1 \;\text{---}\; 0 \;\text{---}\; \infty$, l'arc supérieur marqué $r$, avec l'inscription $\sigma_0 r = \sigma_1 r$.}

\uncertain{Cartes} \qquad $\mathcal{C}_2$-\uncertain{ens}

\uncertain{Cartes orientées} \qquad $\mathcal{C}_2^0$-\uncertain{ens}

\[
\sigma_0^2 = \sigma_1^2 = \sigma_\infty^2 = (\sigma_0\sigma_\infty)^2 = 1
\]
\[
(\rho_0, \rho_1, \rho_\infty \,;\ \rho_\infty \rho_1 \rho_0 = \rho_1^2 = 1)
\]

\[
\begin{array}{lcl}
0 \mapsto 1, & & 1 \mapsto 1, \\
1 \mapsto 1, & -1 \mapsto 0, & 0 \mapsto \infty, \\
i, -i \longmapsto \tfrac{1}{2} & &
\end{array}
\qquad
\varphi(z) = \frac{z-1}{z+1}, \quad \Bigl(\frac{z-1}{z+1}\Bigr)^2 .
\]
\note{Deux symboles biffés ($\sigma$ ? ou un éclair) accompagnent ces correspondances ; $(0, \infty)$ et $1$ sont écrits au-dessous de $z$ dans la même colonne.}

\begin{tikzcd}
X_1 \arrow[r, Rightarrow] & \mathbb{X}_1 \arrow[r, "z \mapsto z^2"] & \mathbb{X}_0 \arrow[r, "\psi"] & \mathbb{X}_0
\end{tikzcd}

\[
(0, 1, \infty) \longrightarrow (0, 1, \infty) \longmapsto 0 \;\; \infty \;\; 1
\qquad\qquad 1 \;\; {-1} \;\; 0 .
\]

\drawing{Une sphère dont l'équateur porte un segment épais de $0$ à $0$, $1$ au milieu, $\infty$ au fond ; une sphère vide avec un arc pointillé vers $0$ ; une grande ellipse avec $\infty$ ; un segment épais de $0$ à $1$.}

\[
\frac{\left(\frac{z-1}{z+1}\right)^2}{\ \cdots\ } = \frac{(z-1)^2}{(z-1)^2 - (z+1)^2}
\]
\note{La dernière ligne est coupée par le bord du feuillet ; le dénominateur de gauche, rendu « $\cdots$ », n'est pas lisible : \ill{}.}

\page{67}
\begin{tikzcd}[column sep=small]
M_{0,5} \arrow[rr, hook] \arrow[dr, hook] & & \hat{M}_{0,5} \\
& \hat{M}'_{0,5} \arrow[ur, hook] &
\end{tikzcd}

$\hat{M}'_{0,5}$ --- correspond aux courbes MD de type $(0,5)$ qui restent stables en \struck{\ill{}} \ill{} \uncertain{l'un des} \ill{} de leurs 5 ;
\[
\hat{M}'_{0,5} \simeq \hat{\mathcal{X}}^*_{0,4} .
\]
\note{Sous $\hat{\mathcal{X}}^*_{0,4}$, un symbole vertical noirci (une égalité ou un isomorphisme biffé).}

\[
\hat{\mathcal{X}}_{0,4} \longrightarrow \mathbb{P}^1 \times \mathbb{P}^1
\]
éclatement des trois points de la diagonale $\Delta$, correspondant aux pts $0, 1, \infty$ ;
\[
\hat{\mathcal{X}}^*_{0,4} \simeq \hat{\mathcal{X}}_{0,4} - \text{les quatre sections } \tilde{\Delta}_0, \tilde{\Delta}_1, \tilde{\Delta}_\infty, \tilde{\Delta} .
\]
\note{Devant le second $\hat{\mathcal{X}}_{0,4}$, un mot biffé : \struck{\ill{}}. Les indices des trois premières sections sont lus avec doute ($\tilde{\Delta}_1$ pourrait être $\tilde{\Delta}_2$).}

\struck{\ill{}} \(\hat{M}_{0,5}\) via \(\hat{\mathcal{X}}^*_{0,4}\) ?

Que manque-t-il : courbes \add{MD} du type $(0,5)$ qui ne sont plus stables en enlevant un des pts : \ill{} \ill{} \ill{} \uncertain{un seul} \ill{} \uncertain{singularité}.

\drawing{Plusieurs configurations de courbes stables : une droite portant deux grappes de composantes, marquée $5$ ; deux droites sécantes dont l'une porte trois points et un point marqué $5^*$, $i$, avec la légende $1 \leqslant i \leqslant 4$ ; une droite portant trois composantes verticales ; trois droites dont l'une porte un point, marquée $\mathbb{P}$.}

\begin{tikzcd}
\hat{\mathcal{X}}_{0,4} \arrow[d] & \hat{\mathcal{X}}_{0,4} \times \hat{\mathcal{X}}_{0,4} \arrow[l] \arrow[d] \\
\hat{M}_{0,4} \arrow[r, no head, "\supset" description] & \hat{\mathcal{X}}_{0,4} \arrow[u, bend right=40, "{\sigma_i \; (i \leqslant 4), \; \delta}"']
\end{tikzcd}

\note{Devant la colonne de gauche un symbole biffé ; $\hat{M}_{0,4}$ est écrit sur un mot biffé. La section notée « $\sigma_i\ (i \leqslant 4)$, $\delta$ » est une flèche courbe remontant de $\hat{\mathcal{X}}_{0,4}$ vers le produit.}

$\hat{M}_{0,3}$ : pts isolés --- \uncertain{en nombre} \ill{}
\[
\hat{M}_{0,3} \;\text{---}\; \hat{\mathcal{X}}_{0,3}
\]

\page{70}
\[
\bigotimes_{i \in K} L_i = L_K
\]
\[
L_i \simeq L_K \otimes p^*_{K \setminus \{i\}, K}\, L_{K \setminus \{i\}}
\]
\note{$L_K$ est entouré et souligné dans la seconde formule ; dessous, un mot en pointillé, \ill{}. Au-dessous encore, une ligne biffée : « $12$ \ill{} $6$ ».}

\[
\begin{array}{l|c|c|c}
n & 4 & 5 & 6 \\ \hline
(n-1)(n-2) & 6 & 12 & 20 \\ \hline
\text{meilleur exposant pour trivialiser } L & 3 & 12\,? & 20\,? \\ \hline
\text{meilleur exposant pour trivialiser } L_i & 6 & 12\,? & 60\,? \\ \hline
\text{exposant pour trivialiser } \text{\struck{$H^1(M_{0,I}, \mathfrak{S}_I;\, \mathbb{G}_m)$}} & & & \\ \hline
\text{dividende (?)} \ H^1(M_{0,I}, \mathfrak{S}_{I \setminus \{i\}};\, \mathbb{G}_m) & & &
\end{array}
\]
\note{« dividende » est une lecture douteuse. Les intitulés des lignes sont écrits obliquement dans la marge gauche du tableau ; les deux dernières lignes sont restées vides. Dans la première colonne, « $(n-1)(n-2)$ » est écrit sur un premier essai noirci.}

$L_K$ \uncertain{quand} $n = 5$ :
\[
\cdots\ \prod_{A \in \mathbb{P}_3}{}' p_A^*(L_A)
= \prod_{\substack{(i \in A) \\ i \in A \in \mathbb{P}_3(K)}} L_i
= \Bigl(\prod_{i \in K} L_i\Bigr)^{\otimes 4}
= L^{4}
\]
\note{Le premier membre, rendu « $\cdots$ », est illisible : \ill{}. Sous $L_A$, un mot en pointillé et « \uncertain{p.\ 3} », illisibles. L'exposant final de $L$ est surchargé.}
\[
(L^4)^3 \simeq \text{\struck{$\mathcal{O}$ \ldots}}
\]

\page{72}
\note{Les notes sont écrites sous un sujet d'examen dactylographié (Université des Sciences et Techniques du Languedoc, session de juin 1978) qui demande de décrire les classes de conjugaison de sous-groupes et d'éléments du groupe alterné $\mathfrak{A}_5$ en termes de l'icosaèdre ; seules les notes manuscrites sont transcrites. La date de l'imprimé est un terme a quo pour cette page.}

\drawing{Trois points alignés ; au-dessous, trois points sur un segment.}

\[
\xi_A \in \Gamma(M_{0,I}, L_A^*) ,
\qquad
\prod_A \xi_A \in \Gamma\Bigl(M_{0,I}, \bigotimes_A L_A\Bigr)
\]
\[
\xi, \quad A \in \mathbb{P}_3(I), \qquad M_{0,I} \xrightarrow{\ \ast\ } M_{0,A}
\]
\[
\xi = \prod \xi_A \in L^{\otimes \nu}, \qquad \tfrac{(n-1)(n-2)}{2} = \nu
\]
\note{L'étiquette $\ast$ de la flèche $M_{0,I} \to M_{0,A}$ est un symbole surchargé illisible. Devant « $\frac{(n-1)(n-2)}{2}$ » un symbole biffé ; le dénominateur est surchargé, et l'exposant de $L$ ($\otimes\nu$) est noirci.}

\[
L_i \xrightarrow{\ \varphi_{i,A,\omega}\ } \mathcal{O}, \qquad i \in A \in \mathbb{P}_3(I),
\qquad
L = \bigotimes_{i \in I} L_i ,
\qquad
\nu = \frac{(n-1)(n-2)}{2}
\]
\[
L_i^2 \xrightarrow[\ \sim\ ]{\ \psi_{i,A}\ } \mathcal{O}
\qquad
\text{\struck{$L_i \xrightarrow{\ \sim\ } \mathcal{O}$}}
\qquad
\xi \in L_i \ \text{---}\ \xi
\]
\[
\xi^*_{i,A} \in \Gamma L_i^2 \qquad \xi \in L_A^2 \simeq
\]

\note{Dans « $L = \bigotimes L_i$ », un symbole biffé précède le produit ; un autre est biffé entre $\xi^*_{i,A}$ et $\xi \in L_A^2$.}

\begin{itemize}
\item \emph{Ex} $n = 4$, $\nu = 3$ : \uncertain{on trouve} $L^{\otimes 6} \simeq 1$ ;
\item \emph{Ex} $n = 5$, $\nu = 6$, \uncertain{on trouve} $L^{\otimes 6} \simeq 1$.
\end{itemize}

\[
L_i \simeq L \otimes L^{-1}_{I \setminus \{i\}} \simeq L \otimes p_i^*\bigl(\mathcal{O}(\ast)\bigr),
\qquad
p_i : M_{0,I} \longrightarrow M_{0, I \setminus \{i\}} \simeq \Sigma
\]
\note{Dans le dernier exposant, l'argument de $\mathcal{O}$ est surchargé (rendu $\ast$ ; peut-être $4$ ou $1$).}

\page{73}
\ill{}\,$\mathcal{T}_{0,I}$ \qquad $\mathcal{T}_{0,I}$ \qquad
\(i_0 \in I\), \quad \(i_1 \neq i_0\), \quad
\(i_0 \in I \setminus \{i_1\} = i_0 \sqcup (I \setminus \{i_0, i_1\})\)

\begin{tikzcd}
\mathcal{T}_{0,I} \arrow[d, "p_1"] & M_{0,K} \arrow[d] \\
\mathcal{T}_{0, I \setminus \{i_1\}} & M_{0,I}
\end{tikzcd}

\note{Devant le premier $\mathcal{T}_{0,I}$ de la colonne, un symbole biffé.}

$\mathcal{T}$ \qquad $\times$ \quad $1\ 2\ 3\ 4$ \quad $5$
\[
\mathcal{O}(4)^{\ast}\ \cdots \simeq \mathcal{O}(-2) \simeq \Omega^1
\]
\[
H \qquad \mathcal{O}( \qquad \prod_{i \in I} L_i
\]
\[
M_{0,I} \hookrightarrow \prod_{i \in K} \Sigma_{J(i)}, \qquad J(i) = I \setminus \{i\}
\]
\note{L'exposant $\ast$ de $\mathcal{O}(4)$ est un petit signe illisible, suivi d'un mot souligné, rendu « $\cdots$ », lui aussi illisible : \ill{}. Dans $\prod L_i$, un symbole est biffé devant $L_i$. Sous $\Sigma_{J(i)}$, souligné, est écrit « $I \setminus \{i\}$ », que l'on rend par l'égalité ; l'indice du produit, $K$, est écrit sur un $I$.}

\[
H^1(\mathfrak{S}_5, M^!_{0,5};\, \mathbb{G}_m), \qquad H^1(M_{0,5};\, \mathbb{G}_m) \simeq 0,
\qquad H^0(M^!_{0,5}, \mathbb{G}_m),
\]
\[
H^2(\mathcal{T}^+_{0,5}, \mathbb{Z}) = \qquad H^1(M^!_{0,5}
\]
\note{Ces lignes sont laissées inachevées ; un trait les sépare de la colonne de droite.}

$M_{0,5}$ à \uncertain{corps} d'\uncertain{éclatement} \ill{} : \uncertain{partir} de $\mathbb{P}_1 \times \mathbb{P}_1$.

\begin{tikzcd}
M^!_{0,5} \arrow[d] & \mathcal{U}_{3} \times \mathcal{U}_{0,3} \arrow[d] \\
M^!_{0,4} \simeq \mathcal{U}_{0,3} & \mathcal{U}_{0,3}
\end{tikzcd}

\[
H^0(M_{0,3}, \mathbb{G}_m) \simeq \text{\struck{$\ast$}}
\]
\note{Le second membre biffé, rendu $\ast$, est illisible.}

\drawing{Deux petites configurations de courbes stables, la première biffée.}

\[
0 \longrightarrow \mathbb{G}_m \longrightarrow H^0(M_{0,K}, \mathbb{G}_m) \longrightarrow \bigl(\mathbb{Z}^{\mathbb{P}_3(K)}\bigr)' \longrightarrow 0
\]
\[
0 \longrightarrow \underbrace{H^1(\mathfrak{S}_K, \mathbb{G}_m)}_{\pm 1}
\longrightarrow \underbrace{H^1\bigl(\mathfrak{S}_K, H^0(M_{0,K}, \mathbb{G}_m)\bigr)}_{H^1(M_{0,K}, \mathfrak{S}_K;\, \mathbb{G}_m)}
\longrightarrow H^1\bigl(\mathfrak{S}_5, (\mathbb{Z}^{\mathbb{P}_3(K)})'\bigr)
\]
\[
H^1\bigl(\mathfrak{S}_5, (\mathbb{Z}^{\mathbb{P}_3(K)})'\bigr) \longrightarrow H^2(\mathfrak{S}_K, \mathbb{G}_m)
\]
\note{Dans la première suite, le $K$ de $M_{0,K}$ et de $\mathbb{P}_3(K)$ est écrit sur un autre indice ; au-dessus de $M_{0,K}$ un petit symbole \ill{}. La dernière flèche est verticale sur la page, de $H^1(\mathfrak{S}_5, \ldots)$ vers $H^2(\mathfrak{S}_K, \mathbb{G}_m)$.}

\page{74}
\note{La page commence au milieu d'un énoncé (« dont la deuxième… ») dont le début ne se trouve pas dans ce lot.}

dont la deuxième est donnée par
\[
\boxed{\rho_1 \longleftrightarrow \sigma_\infty, \quad \rho_\infty \longmapsto \rho^{-1}}
\qquad
\rho_0 = (\rho_\infty \rho_1)^{-1} \longrightarrow \sigma_\infty \rho^{\ast} = \varepsilon_0
\tag{41}
\]
et la première par
\[
\boxed{\tau_0 \longmapsto \tau_\infty}, \quad
\underset{\displaystyle \rho_0}{\tau_0 \,\shortparallel} \longmapsto \quad , \quad
\boxed{\tau_\infty \longmapsto \tau_\infty} ,
\tag{42}
\]
\note{L'exposant $\ast$ de $\rho$ dans (41) est illisible. Dans (42), sous le second $\tau_0$, un signe d'égalité vertical renvoie à $\rho_0$ ; l'image de $\tau_0$ est laissée en blanc. Le $(41)$ est écrit sur un autre numéro. Les indices $0$ de $\tau$ sont noircis et lus avec doute.}

\drawing{Un triangle sphérique de sommets $0$, $1$, $\infty$, avec un signe $\xi$ (ou $\zeta$) à gauche et un $i$ au-dessous ; à droite, une figure mêlant un triangle (sommets marqués, dont l'un $\infty$ biffé et l'autre $1$) et un pentagone, une région hachurée en zigzag.}

\[
\pi_{\mathbb{D}_n}, \quad \pi_{\mathfrak{S}_4^+}, \quad \pi_{\mathfrak{S}_4}, \quad \pi_{\mathfrak{S}_5^+}
\]

$\pi_X$ \qquad $\pi_X \hookrightarrow X$ une carte : \uncertain{on enlève} les \uncertain{sommets}

$\pi_{P}$\ill{} \qquad $\pi_{\mathrm{icos}}$, $\pi_{\mathrm{dodéc}}$, $\pi_{\mathrm{tétr.}}$, $\pi_{\mathrm{cube}}$, $\pi_{\mathrm{oct.}}$, \struck{$\pi$\ill{}}

\qquad $\pi$\ill{}, \qquad $\pi_{\mathbf{P}} \longrightarrow$ polygone combinatoire fini…

\page{76}
\drawing{Feuillet en largeur, presque entièrement de figures : trois sphères à gauche, dont l'équateur porte des points cerclés étiquetés $a, b, c, d, e$ et $s, t, u, v, w$, avec un pôle $a$ relié par des méridiens, un secteur hachuré et des flèches de rotation ; au centre, un grand patron (un pentagone en perspective, sommets $s$, $t$, et des faces hachurées) dont les arêtes portent $(a), (b), (c), (d), (e)$ ; au-dessous, une sphère et deux cercles équatoriaux portant les mêmes lettres ; en haut à droite, un fuseau de sommets $s$ et $t$, arêtes $b, c, d, e$, une région hachurée et un point $\alpha$.}

\begin{tikzcd}[row sep=small, column sep=large]
& \mathfrak{F} \\
\mathcal{R} \arrow[ur] \arrow[r] \arrow[dr] & \mathcal{A} \\
& \mathcal{S}
\end{tikzcd}

\begin{itemize}
\item $\mathcal{R}$ : repères des \emph{bi}dodécagones (120) \uncertain{grands} ;
\item $\mathfrak{F}$ : faces des bidodécagones (12) ;
\item $\mathcal{A}$ : arêtes \struck{des} \struck{\ill{}} des dodécagones (30) ;
\item $\mathcal{S}$ : sommets des bidodécagones (20).
\end{itemize}
\note{Le nombre des faces est surchargé : « 12 » est lu avec doute. Le $\mathfrak{F}$ est écrit sur une autre lettre.}

\[
2 \times (6 + 15 + 10) = 62
\]

\page{77}
\note{Croquis sans texte de sa main, dessinés sur une circulaire datée du 12 février 1982 (terme a quo pour cette page). Ils prolongent les figures de cartes et de recollements des pages précédentes.}

\drawing{Un ruban hexagonal hachuré, replié sur lui-même ; plusieurs graphes trivalents en forme de « nœud papillon » ou de deux triangles joints par une arête, certains avec des sommets cerclés ou hachurés ; une grande ellipse au crayon traversée par un segment ; un réseau en zigzag de quadrilatères et de triangles ; un segment portant trois points ; une boucle à deux points cerclés ; un triangle dont un sommet est cerclé, prolongé par une queue.}

\page{78}
\note{Feuillet de chemise, vide à part un titre de sa main en haut à droite (et, au crayon, l'indication « (19\,p) », peut-être des archivistes). Il marque le début d'une nouvelle liasse, qui commence à la page 80.}

Groupoïdes

\page{80}
\emph{Lemme 1.} Soit
\[
1 \to L \longrightarrow \tilde{\pi} \longrightarrow \pi \longrightarrow 1
\]
une suite exacte de groupes, soit $(x_i)_{i \in I}$ un syst.\ de générateurs de $\pi$, $(a_j)_{j \in J}$ \uncertain{un syst.\ de générateurs de} $L$,
\[
R_\alpha\bigl((x_i)\bigr) = 1 \quad (\alpha \in A)
\]
un syst.\ de relations définissantes de $\pi$,
\[
S_\beta\bigl((a_j)\bigr) = 1 \quad (\beta \in B)
\]
un syst.\ de relations définissantes de $L$. Relevons les $x_i$ en des $\tilde{x}_i \in \tilde{\pi}$, et choisissons des mots $M_\alpha$ et $M_{ij}$ ($\alpha \in A$, \struck{\ill{}} $(i,j) \in I \times J$) en des lettres $a_j$, de telle façon qu'on ait
\[
\left\{
\begin{aligned}
& R_\alpha\bigl((\tilde{x}_i)\bigr) = M_\alpha\bigl((a_j)\bigr) && \forall \alpha \in A \qquad (1) \\
& \tilde{x}_{i_0}\, a_{j_0}\, \tilde{x}_{i_0}^{-1} = M_{i_0 j_0}\bigl((a_j)\bigr) && \forall (i_0, j_0) \in I \times J \qquad (2)
\end{aligned}
\right.
\]
Alors les relations (1), (2), jointes à
\[
S_\beta\bigl((a_j)\bigr) = 1 \quad (\beta \in B) \qquad (3)
\]
sont un syst.\ de relations définissantes (pour $\tilde{\pi}$) par rapport aux générateurs $(\tilde{x}_i)$, $(a_j)$.

\marginal{(1) : relations « \uncertain{relèvements} » des $R_\alpha$ ; (2) : relations \uncertain{d'action} \ill{} ; (3) : « relations de $S_\beta$ » dans $L$.}
\note{Les trois annotations marginales sont écrites en oblique dans la marge gauche, face aux formules (1)--(3) ; leur lecture est très incertaine, en particulier celle qui concerne (2).}

\emph{NB} Si $L$ est central dans $\tilde{\pi}$, \struck{les relations} \add{\uncertain{(relations triviales)}}, on peut prendre les $M_{i_0 j_0}$ les mots vides, donc (2) devient
\[
\tilde{x}_i\, a_j\, \tilde{x}_i^{-1} = a_j, \quad \text{i.e.}\ [\tilde{x}_i, a_j] = 1 \qquad (2')
\]
pour tout $(i,j) \in I \times J$.

\note{La page s'arrête ici ; la suite du lemme et sa démonstration, s'il y en a une, sont hors de ce lot.}

\end{document}
